\textbf{Proof.}
Write \(m=m_n\).  Since \(0<\rho<1\) and \(0<\gamma\leq1\),
\[
0\leq m=\lfloor\rho n^{\gamma}\rfloor< n,
\tag{1}
\]
which proves the stated finite range.  If \(m=0\), then \(G_n\) is the edgeless graph and the all-singleton partition gives
\[
n\left(\frac1n\right)^2=\frac1n
=\frac{m^2+n-m}{n^2}.
\tag{2}
\]
If \(m\geq1\), use the \(m\) clique vertices as one partition block and every other vertex as a singleton.  Applying \(\mathrm{prop{:}clique\text{-}partition\text{-}rademacher\text{-}converse}\) gives
\[
R_2(G_n)\geq
\left(\frac mn\right)^2+(n-m)\left(\frac1n\right)^2
=\frac{m^2+n-m}{n^2}.
\tag{3}
\]
Equations (2)--(3) prove the certificate at every finite \(n\), including the floor boundary \(m=0\).

We next verify the graph quantities before taking orders.  If \(m=0\), all \(n\) vertices are isolated, so \(\alpha(G_n)=n\), \(d_{\mathrm{avg}}(G_n)=0\), and \(\lambda(G_n)=0\).  If \(m\geq1\), an independent set contains at most one clique vertex and can contain all \(n-m\) isolated vertices, whence
\[
\alpha(G_n)=n-m+1.
\tag{4}
\]
The graph has exactly \(\binom m2\) edges, so the degree-sum identity gives
\[
d_{\mathrm{avg}}(G_n)=\frac{2\binom m2}{n}=\frac{m(m-1)}n.
\tag{5}
\]
Finally, its adjacency matrix is block diagonal with one \(K_m\) block and zero singleton blocks.  The largest eigenvalue of the \(K_m\) block is \(m-1\), proving the stated formula for \(\lambda(G_n)\), including \(m=1\).

For the asymptotics, \(\rho n^{\gamma}\to\infty\), so for all sufficiently large \(n\),
\[
\frac{\rho}{2}n^{\gamma}\leq m\leq\rho n^{\gamma}
\quad\text{and}\quad m\geq2.
\tag{6}
\]
Also \(n-m=\Theta(n)\): if \(\gamma<1\), then \(m/n\to0\), while if \(\gamma=1\), then \(n-m\geq(1-\rho)n\).  Hence (3) and (6) imply
\[
\frac{m^2}{n^2}=\Theta(n^{2\gamma-2}),
\qquad
\frac{n-m}{n^2}=\Theta(n^{-1}).
\tag{7}
\]
The sum of two nonnegative sequences is of the order of their maximum, proving
\[
L_n=\Theta\!\left(n^{-1}\vee n^{2\gamma-2}\right).
\tag{8}
\]

Equations (4)--(6) similarly give
\[
\frac1{\alpha(G_n)}=\Theta(n^{-1}),
\qquad
\frac{\sqrt{d_{\mathrm{avg}}(G_n)}}n
=\frac{\sqrt{m(m-1)}}{n^{3/2}}
=\Theta(n^{\gamma-3/2}),
\tag{9}
\]
and, because \(m\geq1\),
\[
\frac{\lambda(G_n)+1}{n}=\frac mn=\Theta(n^{\gamma-1}).
\tag{10}
\]
Applying \(\mathrm{lem{:}published\text{-}dte\text{-}graph\text{-}functional\text{-}bounds}\) after these exact calculations proves the two asserted published-functional orders.

If \(\gamma\leq1/2\), both lower displays are dominated by \(n^{-1}\).  If \(\gamma>1/2\), their dominant terms are respectively \(n^{2\gamma-2}\) and \(n^{\gamma-3/2}\), whose ratio is \(n^{\gamma-1/2}\).  Finally, equality of the clique-lower and spectral-upper exponents requires \(2\gamma-2=\gamma-1\), hence \(\gamma=1\).  At the boundary \(\gamma=1/2\), both lower orders are exactly \(n^{-1}\); at \(\gamma=1\), the clique lower and spectral upper are both order one.  These boundary checks establish the stated elbow and matching verdict. \(\square\)